\textbf{Figure 1. A harmonized multimodal human MASLD Resource.}

\textbf{(A)} Analysis design. Human liver samples (healthy, MASL, MASH) enter three integrated arms: genetics (genome-wide association, fine-mapping, colocalization with liver eQTL), transcriptomics (uniform alignment, differential expression, gene programs, cell--cell communication) and molecular context assays (chromatin accessibility, liver protein, spatial transcriptomics, clinical drug databases). Each gene's evidence across arms is recorded in a gene-by-evidence matrix, the MASLD Gene Catalog (Figure 6). Arrows show analysis order, not causal links.

\textbf{(B)} Datasets by modality (human cohorts only). Inner ring: modality and number of datasets; middle ring: condition (disease, control) or GWAS ancestry (EUR, AFR, EAS, AMR, SAS); outer ring: species; bars: samples per dataset. Totals per modality are printed beside each segment (see Methods for updated deduplicated donor and library counts). The inferential unit differs by arm (Methods): bulk RNA uses genetic-identity clusters, single-cell analyses use donors, and cells, spots and runs do not add to the biological sample size.
